%% =====================================================================
%%  10 Appendix
%% =====================================================================
%%  GUIDE
%%   - Long algebra, code, raw data tables, extra parameter studies.
%%   - The method must still be explained in the body (official rule);
%%     the appendix is for details, not for hiding the method.
%%   - Code is quoted from code/ so the listing and the run stay identical.
%% =====================================================================
\appendix

\section{Derivation of the orifice speed}
\label{app:derivation}

Starting from the steady energy equation for an adiabatic nozzle with no
work, $h+u^2/2=h_0$, and using $h=c_pT$ with
$T/T_0=(p/p_0)^{(\gamma-1)/\gamma}$, the exit speed is
\begin{equation}
  u^2 = \frac{2\gamma}{\gamma-1}\frac{p_0}{\rho_0}
  \left[1-\left(\frac{p_a}{p}\right)^{(\gamma-1)/\gamma}\right],
  \label{eq:app-derivation}
\end{equation}
which is Eq.~\eqref{eq:subsonic} written at the reservoir state. The
bracket is largest when the flow chokes, at
$p/p_a=((\gamma+1)/2)^{\gamma/(\gamma-1)}$; substituting that ratio gives
\begin{equation}
  u_{\rm ch}^2 = \frac{2\gamma}{\gamma+1}\frac{p_0}{\rho_0}
  = \frac{2\gamma}{\gamma+1}RT_0,
  \label{eq:app-choked}
\end{equation}
that is, Eq.~\eqref{eq:choked}, whose speed no longer depends on the
ambient pressure. Differentiating the mass balance
$\dd m/\dd t=-\rho_e A u$ with $m=\rho V$ and using
$\rho\propto p^{1/\gamma}$ gives Eq.~\eqref{eq:mass}, with the ratio
$\rho_e/\rho=(2/(\gamma+1))^{1/(\gamma-1)}=0.63$ absorbed into $C_d$ ---
which is why the fitted coefficient is not the handbook one.

\section{Numerical code}
\label{app:code}

The integration quoted in Section~\ref{sec:method} is performed by the
script below; the listing is included with \texttt{lstinputlisting} so the
printed code cannot drift from the version that produced the numbers.

\lstinputlisting[language=Python,caption={Blow-down integration used for
all results in Section~\ref{sec:results}.}]{code/blowdown.py}

\section{Parameter sweep data}
\label{app:data}

\begin{table}[htbp]
  \centering
  \small
  \caption{The three diameters used in Table~\ref{tab:results}. The last
  column multiplies the measured time by $(d/\SI{10}{cm})^{1.97}$; it is
  constant to \SI{1}{\percent}, which is what makes the power law a valid
  summary over this range.}
  \label{tab:sweep}
  \begin{tabular}{@{}S[table-format=2.0]S[table-format=4.1]S[table-format=3.1]@{}}
    \toprule
    {$d$ (cm)} & {$t_{0.3}$ (s)} & {$t_{0.3}(d/\SI{10}{cm})^{1.97}$ (s)} \\
    \midrule
     5 & 1108 & 283 \\
    10 &  286 & 286 \\
    20 &   73 & 284 \\
    \bottomrule
  \end{tabular}
\end{table}
